% TOP_PROBLEM: {"id": 1457, "title": "Rota's Basis Conjecture", "category_id": 4, "category": {"id": 4, "name": "algebra", "display_name": "Algebra", "description": "Group theory, ring theory, field theory, and algebraic structures.", "slug": "algebra", "order_index": 4, "created_at": "2026-07-31T15:26:25.671Z"}, "status": "open"}
% FIRST_POSTED: 2026-10-01
% ATTEMPT_STATUS: unresolved
% Imported from: unsolvedmath_additional_500.tex; UnsolvedMath ID 1457
% !TeX program = lualatex
% Compile twice: lualatex 1457.tex
% Self-contained: no external bibliography, images, or \input files.
% Numeric section labels and visible numbers are original UnsolvedMath IDs.
\documentclass[11pt,a4paper]{article}
\usepackage[margin=24mm,headheight=14pt]{geometry}
\usepackage{fontspec}
\setmainfont{Latin Modern Roman}
\setsansfont{Latin Modern Sans}
\setmonofont{Latin Modern Mono}
\usepackage{amsmath,amssymb,mathtools}
\usepackage{unicode-math}
\setmathfont{Latin Modern Math}
\usepackage{microtype}
\usepackage{enumitem,xurl,needspace}
\usepackage[dvipsnames]{xcolor}
\usepackage{hyperref}
\hypersetup{unicode=true,colorlinks=true,linkcolor=MidnightBlue,urlcolor=MidnightBlue,
 pdftitle={UnsolvedMath Problem 1457},pdfauthor={Source-annotated compilation},bookmarksnumbered=true}
\usepackage{bookmark,tocloft,titlesec,fancyhdr}
\setlength{\cftsecnumwidth}{4.2em}
\cftsetpnumwidth{2.5em}
\cftsetrmarg{4em}
\titleformat{\section}{\Large\bfseries\raggedright}{\thesection}{0.7em}{}
\pagestyle{fancy}\fancyhf{}
\fancyhead[L]{\small\sffamily Additional UnsolvedMath statements}
\fancyhead[R]{\small Original catalogue IDs}
\fancyfoot[C]{\thepage}
\setlength{\parindent}{0pt}
\setlength{\parskip}{5pt plus 1pt}
\setlength{\emergencystretch}{3em}
\setcounter{tocdepth}{1}\setcounter{secnumdepth}{1}
\setlist[itemize]{leftmargin=3em,itemsep=4pt,topsep=4pt,parsep=0pt}
\allowdisplaybreaks
\newcommand{\N}{\mathbb{N}}
\newcommand{\Z}{\mathbb{Z}}
\newcommand{\Q}{\mathbb{Q}}
\newcommand{\R}{\mathbb{R}}
\newcommand{\C}{\mathbb{C}}
\newcommand{\F}{\mathbb{F}}
\newcommand{\A}{\mathbb{A}}
\newcommand{\PP}{\mathbb{P}}
\newcommand{\E}{\mathbb{E}}
\newcommand{\eps}{\varepsilon}
\newcommand{\dd}{\,\mathrm d}
\newcommand{\sref}[2]{\hyperlink{source:#1:#2}{[S#2]}}
\newif\ifshowcatalogue
\showcataloguetrue
\newcommand{\cataloguescope}[1]{\ifshowcatalogue
 \paragraph{Statement recorded in the supplied source.}
 #1\fi}
\begin{document}
% UnsolvedMath ID 1457; source catalogue code ALG-012

\setcounter{section}{1456}

\section{Rota’s Rainbow-Basis Conjecture}\label{1457}

{\small\sffamily UnsolvedMath ID 1457 · Algebra}

\subsection{Definitions and mathematical statement}

\paragraph{Shared notation.}Unless a section says otherwise, $\N=\{1,2,\ldots\}$,
$\Z$ is the ring of integers, $\Q,\R,\C$ are the rational, real, and complex
fields, $[N]=\{1,\ldots,\lfloor N\rfloor\}$, and $\log$ is the natural
logarithm. For real functions of a parameter tending to infinity,
$f\ll g$ and $f=O(g)$ mean $|f|\leq Cg$ eventually for some fixed $C$;
$f\gg g$ reverses this comparison; $f=o(g)$ means $f/g\to0$;
$f\sim g$ means $f/g\to1$; and $f\asymp g$ means both order bounds.
A subscript on $O$ or $\ll$ specifies allowed constant dependence.
The expression $n^{o(1)}$ in an upper bound means growth smaller than
$n^\epsilon$ for each fixed $\epsilon>0$. Local definitions govern any
exception, finite-field convention, or use of the same symbol for another
object. An asymptotic claim is not automatically an effective algorithm.



A matroid of rank $n$ has a collection of independent subsets satisfying heredity and augmentation, with every basis of size $n$. Regard the supplied $n$ bases as disjoint labelled copies when an element appears in more than one. Seek a partition of their $n^2$ occurrences into $n$ bases, each taking exactly one occurrence from each original basis. This is a rainbow decomposition, not merely one rainbow basis. \sref{1457}{1} \sref{1457}{2}

\paragraph{Statement recorded in the supplied source.}
Given n bases of an n-dimensional matroid, can we find n disjoint rainbow bases? \sref{1457}{1}

\paragraph{Residual task reported by the archive.}

The embedded review (2026-08-17) records: Prove the rainbow-basis decomposition universally or find a counterexample. \sref{1457}{1}

\subsection{Short English statement}

Can n bases of a rank-n matroid be rearranged into n new bases, each containing one element from every original basis?

\subsection{Research attempt: the rank-two case, a partition-matroid construction, and a greedy obstruction}
\textbf{Status and source correction.} The conjecture is not proved in general. The inspected arXiv record [S2] lists Alexey Pokrovskiy as its sole author; its abstract proves an asymptotic statement with both the number and sizes of the independent sets equal to $n-o(n)$, not an exact decomposition. We prove the following exact subcases independently.

\subsubsection{1. Two bases of a rank-two matroid}
Write the bases as labelled occurrences $\{a,b\}$ and $\{c,d\}$. None is a loop. Among nonloops in a rank-two matroid, dependence of a pair is exactly membership in the same parallel class, and parallelism is an equivalence relation. If both $\{a,c\}$ and $\{b,d\}$ are independent, they are the desired rainbow bases. Otherwise, after interchanging names within the rows if needed, assume $a,c$ are parallel. Since $c,d$ is independent, $a,d$ is independent. Since $a,b$ is independent, $b,c$ is independent. Therefore the other pairing $\{a,d\},\{b,c\}$ works. This also handles repeated underlying elements, because their separate occurrences are parallel.

\subsubsection{2. A cyclic construction for partition matroids}
Consider a rank-$n$ partition matroid with $n$ nonempty classes and capacity one in each class. Every basis contains one element of each class. Label rows by $i\in\mathbb Z/n\mathbb Z$, and within row $i$ denote its occurrence from class $j$ by $b_{ij}$. For column $k$, take
\[C_k=\{b_{i,i+k}:i\in\mathbb Z/n\mathbb Z\}.
\]
Each column has one occurrence from each row. Its class labels $i+k$ run through all classes, so it is a basis. For fixed $i$, as $k$ varies every $b_{ij}$ appears exactly once; hence the columns partition all occurrences. This proves the exact conjecture for every such partition matroid, including the repeated-identical-basis configuration.

\subsubsection{3. One good transversal need not leave a solvable remainder}
Already over $\mathbb Q^2$, take rows
\[B_1=\{e_1,e_2\},\qquad B_2=\{e_1+e_2,e_2\}.
\]
Choosing $\{e_1,e_1+e_2\}$ gives a perfectly valid rainbow basis, but the remaining pair consists of two labelled copies of $e_2$ and is dependent. The alternative first choice $\{e_1,e_2\}$ using the second-row $e_2$ leaves $\{e_2,e_1+e_2\}$, also a basis. Thus an existence theorem for a single rainbow basis cannot simply be iterated with arbitrary choices.

The obstruction is structural: after removing one transversal, each row has $n-1$ elements but need not be a basis of one common rank-$(n-1)$ matroid obtained by a shared contraction. Contracting different selected elements separately in different rows does not define a single residual matroid on all occurrences. A proof using induction must preserve a common independence structure, not just the right number of elements in each row.

\subsubsection{4. Verification and exact remaining gap}
The diagnostic checks the two rank-two pairings for all pairs of bases represented by nonzero vectors in $\mathbb F_p^2$ for $p=2,3,5$, computes the displayed rational determinants, and checks every row and class count in the cyclic construction through $n=20$. The rank-two proof rests on parallel-class equivalence, which does not characterize dependence in higher rank: three vectors can be dependent while every pair is independent. The partition construction similarly assumes independence depends only on class multiplicities. A mechanism handling higher-order dependencies while choosing all $n$ transversals simultaneously remains missing. The verified subcases and the greedy counterexample are partial work, not an exact solution inferred from an asymptotic theorem.

\subsection{Sources}

{\small

\begin{itemize}

\item[\hypertarget{source:1457:1}{S1}] ULAM.AI, \emph{UnsolvedMath}, supplied \texttt{problems.json}, record 1457 (ALG-012), statement and background. The embedded literature-review date is 2026-08-17; that is the archive’s review date, not a fresh certification. Dataset landing page: \url{https://huggingface.co/datasets/ulamai/UnsolvedMath}.

\item[\hypertarget{source:1457:2}{S2}] A. Pokrovskiy, Rota's Basis Conjecture holds asymptotically, arXiv:2008.06045. \url{https://arxiv.org/abs/2008.06045}. Bibliographic information imported from the supplied record.

\end{itemize}

}

\label{end:1457}
\end{document}
